% =====================================================================
%  Comandi e ambienti
% =====================================================================

% --- Enunciati -------------------------------------------------------------
\theoremstyle{definition}
\newtheorem{definizione}{Definizione}[chapter]
\newtheorem{ipotesi}{Ipotesi}
\theoremstyle{plain}
\newtheorem{proposizione}[definizione]{Proposizione}
\theoremstyle{remark}
\newtheorem{osservazione}[definizione]{Osservazione}
\renewcommand{\qedsymbol}{$\blacksquare$}
\crefname{definizione}{definizione}{definizioni}
\Crefname{definizione}{Definizione}{Definizioni}
\crefname{proposizione}{proposizione}{proposizioni}
\Crefname{proposizione}{Proposizione}{Proposizioni}
\crefname{osservazione}{osservazione}{osservazioni}
\Crefname{osservazione}{Osservazione}{Osservazioni}

% --- Notazione matematica ----------------------------------------------------
\newcommand{\R}{\mathbb{R}}
\newcommand{\N}{\mathbb{N}}
\newcommand{\E}{\mathbb{E}}
\newcommand{\Ind}{\mathbf{1}}
\DeclareMathOperator*{\argmin}{arg\,min}
\DeclareMathOperator*{\argmax}{arg\,max}
\DeclareMathOperator{\softmax}{softmax}
\DeclareMathOperator{\sg}{sg}                 % stop-gradient
\DeclareMathOperator{\GAP}{GAP}               % global average pooling
\DeclareMathOperator{\PPL}{PPL}               % perplexity del codebook
\DeclareMathOperator{\KL}{D_{KL}}
\newcommand{\Hent}{\mathrm{H}}
\newcommand{\norm}[1]{\left\lVert #1 \right\rVert}
\newcommand{\abs}[1]{\left\lvert #1 \right\rvert}
\newcommand{\T}{^{\top}}
\newcommand{\bigO}[1]{\mathcal{O}\!\left(#1\right)}
\newcommand{\vect}[1]{\bm{#1}}
\newcommand{\hatq}{\hat{\vect{q}}}
\newcommand{\hatk}{\hat{\vect{k}}}

% --- Nomi ricorrenti -----------------------------------------------------------
\newcommand{\maresunet}{MAResU-Net}
\newcommand{\maresunetvq}{MAResU-Net-VQ}
\newcommand{\rsseg}{\texttt{rsseg}}
\newcommand{\RQ}[1]{\textbf{RQ#1}}
\newcommand{\codice}[1]{\texttt{#1}}
\newcommand{\tile}[1]{\texttt{#1}}

% --- Quadratino con il colore di una classe ISPRS (per le legende) -------------
\newcommand{\ClasseBox}[1]{%
  \tikz[baseline=-0.6ex]\draw[black!60, line width=0.3pt, fill=#1] (0,0) rectangle (0.8em,0.8em) ++(0,-0.4em);}
\newcommand{\ClasseBoxX}[1]{%
  \tikz[baseline=0.05em]\draw[black!60, line width=0.3pt, fill=#1] (0,0) rectangle (0.8em,0.8em);}

% --- Segnaposto per le parti sperimentali ------------------------------------
% \Segnaposto{...}  nota in linea, ben visibile, per contenuti ancora da produrre
% \ND               valore non ancora disponibile in una cella di tabella
% \FiguraMancante{descrizione}{altezza}  riquadro al posto di una figura
\newcommand{\Segnaposto}[1]{\todo[inline]{#1}}
\newcommand{\ND}{\textcolor{SapRosso}{\textit{n.d.}}}
\newcommand{\FiguraMancante}[2][4cm]{%
  \ifbozza
    \begin{tikzpicture}
      \draw[SapRosso, dashed, line width=0.8pt, fill=SapRossoVelo]
        (0,0) rectangle (0.95\linewidth, #1);
      \node[text width=0.8\linewidth, align=center, font=\small\itshape, text=SapRossoScuro]
        at (0.475\linewidth, #1/2) {Figura da inserire\\[0.3em]#2};
    \end{tikzpicture}%
  \else
    \fbox{\parbox[c][#1][c]{0.95\linewidth}{\centering\small\itshape Figura da inserire: #2}}%
  \fi}

% --- Stili TikZ condivisi dagli schemi -------------------------------------
\tikzset{
  >={Stealth[length=2.2mm]},
  blocco/.style={draw=SapRossoScuro, line width=0.5pt, rounded corners=1.5pt,
                 align=center, font=\footnotesize, minimum height=7mm, inner sep=2.5pt},
  enc/.style={blocco, fill=ColEncoder},
  dec/.style={blocco, fill=ColDecoder, draw=SapGrigio},
  attn/.style={blocco, fill=ColAttn},
  quant/.style={blocco, fill=ColQuant, text=white, draw=SapRossoScuro},
  dato/.style={blocco, draw=SapGrigio, fill=white, rounded corners=0pt},
  nota/.style={font=\scriptsize\itshape, text=SapGrigio, align=center},
  flusso/.style={->, draw=ColFreccia, line width=0.6pt},
  grad/.style={->, draw=ColGradiente, line width=0.6pt, dashed},
  skipc/.style={->, draw=SapRosso, line width=0.6pt},
  gruppo/.style={draw=SapGrigioTenue, dashed, rounded corners=3pt, inner sep=6pt},
}
